\subsection{The actual marked ring and propagation through coefficient curves}
\label{sec:selmer-marked-propagation}

The local prime-inclusion theorem and the patched-support theorem now
apply to their prescribed arithmetic rings.  This supplies the full
marked propagation contract, with every prime link and the original
integral equation ideal retained~\cite{SelmerMarkedSource}.

\begin{theorem}\label{thm:selmer-marked-propagation}
Let $\ell\ne2$, and retain the original ordinary determinant pair
$u=(D,\psi)$ over $B$, coefficient DVR $O$, adapted Taylor--Wiles datum
$T$, weight rank $D_{\rm wt}$, and actual modules with base ideal $J$.
Assume $B$ is topologically generated over $O$ by the pair values.
Let
\[
 f_i:B\longrightarrow A_i=k_i[[t_i]],\qquad 0\le i\le h,
\]
be any finite family of the original normalized points, with the
original odd Selmer families, distinct auxiliary eigenvalues, and
bounds B1/B2 with $D_0=D_{\rm wt}+nq$.
Assume the finite $k_i$ have the same cardinality.  Put $p_i=\ker f_i$.
Let $e_i\subseteq p_i,p_{i+1}$ be primes at which the original ordinary
Cayley--Hamilton equation algebra is nonsemisimple after a residue-field
extension.  If
\[
 q_0\subseteq p_0,\qquad J\subseteq q_0,\qquad
                         \operatorname{DiscSt}(D,q_0),
\]
for a prime $q_0$, then every prime $q_*\subseteq p_h$ which is stable
or has the same nonsemisimplicity property contains $J$.
Here stability is exactly the original predicate
\[
 \operatorname{DiscSt}(D,p):\quad
 \det\bigl(\operatorname{tr}D(g_a h_b)\bigr)_{a,b=1}^{n^2}\notin p
                  \quad\hbox{for some }g_a,h_b.             \tag{MP1}
\]
The conclusion holds for every finite chain, with its original kernels
and linking primes.  It is the complete source contract \texttt{Raw},
including its \texttt{Core} and \texttt{CoreBC} forms.
\end{theorem}

\begin{proof}
We give the actual construction, so that no patched-support or
marked-model existence hypothesis is hidden in the statement.

\emph{A single limiting group.}
First take the maximal pro-$\ell$ quotients of the fixed open residual
kernels; their harmlessness is proved below.  The uniform bound on
$\operatorname{Hom}(H_N,\mathbb F_\ell)$ gives a uniform number of
pro-$\ell$ generators.  A free profinite group $P$ on these generators,
fixed lifts of the finite base quotient, and the original conjugation
therefore surjects onto every $\Gamma_N$, compatibly modulo the fixed
open kernel.  Fix a free ultrafilter.  In every finite quotient $P/V$,
the images of $\ker\pi_N$ stabilize.  Compactness gives a closed normal
$K\subset P$ with
\[
                    (\ker\pi_N)V=KV
                  \quad\hbox{for almost all }N.             \tag{MP2}
\]
Put $\Gamma_\infty=P/K$.  Finitely many homomorphisms into each finite
group give compatible limits of the base and local maps, uniformly on
the common cover.  For each local group $Y$, its maps factor through the fixed finitely
generated quotient $Y/\operatorname{PL}(f_0^{-1}H)$, by local
class-field finiteness.  Choose a free profinite presentation of this
quotient.  Lift its finitely many generator images individually to $P$
at every level, and take their ultralimits.  Every relator maps into
$K$ by (MP2), so the limit descends to the local quotient.  Congruence
on the finite generator set gives the required uniform congruence on
the whole group.  A homomorphic lift of $Y$ to $P$ was never assumed.
The limiting auxiliary
Frobenius and inertia commute: in each finite pro-$\ell$ quotient,
$Q_N\equiv1\pmod{\ell^N}$ makes the tame commutator trivial eventually.
Conjugation and every original local subgroup are preserved.  This
constructs the full \texttt{Lim}; its index-two equation group $P^+$
is topologically finitely generated.  Every limiting eventuality below
uses this same ultrafilter, exactly as in Lim and S5; the only use of
cofinite bounds is that they belong to it.

\emph{The marked coefficient ring.}
First suppose $O\twoheadrightarrow k_i$ and that one original curve
variable $s$ satisfies $f_i(s)=t_i$.  Put
\[
 \Lambda=O[[s,X_1,\ldots,X_{D_{\rm wt}}]],\qquad
 S_\infty=\Lambda[[Y_{ji}]]_{1\le j\le q,\,0\le i<n}.
\]
The continuous determinant of $P^+$ and its original residual value
have a complete Noetherian universal ring.  Chenevier's representability
and Wang--Erickson's finite-$H^1$ criterion apply because $P^+$ is
finitely generated~\cite{SelmerMarkedChenevier,SelmerMarkedWE};
there is no requirement $\ell>n$.
Take its completed tensor product $R_1$ with the finitely many local
character deformation rings, retaining the prescribed values on
$\mu_v$ and their finite torsion relations.  Define
\[
 R=R_1[[s,Z_{ji},Y_{ji}]],\qquad U_{ji}=u^0_{ji}+Z_{ji}.     \tag{MP3}
\]
This represents exactly the original determinant, ordered characters,
curve variable, auxiliary roots and diamond variables before imposing
the global equations.  It is the source's \texttt{RRObj}.

Its universal property gives maps $\rho_N:R\to R_N$ into the actual
auxiliary coefficient rings.  The original map $S_\infty\to R$ is
\[
 s\longmapsto s,\qquad
 X_a\longmapsto c_a\psi_{v_a i_a}(\gamma_a)-1,\qquad
 Y_{ji}\longmapsto Y_{ji}.                                  \tag{MP4}
\]
The residual weight identity puts these values in $\mathfrak m_R$;
completeness defines the map, and equality on the variables proves its
compatibility with every original module action.

Let $C/B$ be the local splitting algebra of the $q$ original
polynomials $\chi_D(\sigma_j)$, with their prescribed residual ordered
roots.  The full splitting algebra is finite free of rank $(n!)^q$.
Its selected local factor is a nonzero direct summand over complete
local $B$, hence finite free and faithfully flat.  Its actual rank is
retained.  The universal map
\[
 \pi:R\longrightarrow C,\qquad U_{ji}\longmapsto u_{ji},
                         \quad Y_{ji}\longmapsto0            \tag{MP5}
\]
is surjective: its compact image contains all pair coefficients, $s$,
and the roots, which topologically generate $C$.

\emph{The exact equation quotient.}
Let $I\subset R$ be generated by the original six scalar relation
families: traces of the limiting-group kernel, every polarization
coefficient, the ordered ordinary products, the prescribed tame
inertia, the auxiliary characteristic polynomial, and its partial
polynomial times $\tau_j-1-Y_{ji}$.  Thus the last two are
\[
 \operatorname{tr}D\!\left(g\prod_i(\phi_j-U_{ji})\right),\qquad
 \operatorname{tr}D\!\left(
 g\prod_{a\ne i}(\phi_j-U_{ja})(\tau_j-1-Y_{ji})\right).
                                                               \tag{MP6}
\]
All test elements and original local lifts are quantified.  This is
literally $I=(\texttt{Rel})$, before any localization or radical.
It is closed because $R$ is complete Noetherian.  Every relation maps
to zero under $\pi$, so $\bar\pi:R/I\twoheadrightarrow C$.

The full determinant descends to the limiting group.  Indeed the finite
continuous Cayley--Hamilton algebra has a generalized matrix algebra
with rank-one residual corners.  A two-sided trace-radical ideal has
zero diagonal corners.  Every term involving it in the canonical
cycle formula for the determinant vanishes, after every scalar
extension.  Thus the kernel trace equations kill the full determinant
law without dividing by $n!$.  The remaining equations and their
polarized partners give the prescribed pair and auxiliary relations.
Its equation algebra uses precisely the already closed ideal
\[
             \mathcal I=D.\operatorname{chBar}
                  (\texttt{pairGens}\cup\texttt{qGens}).     \tag{MP7}
\]
The algebraic span of its image is used after base change, as in the
original contract.

Lift the entire prime chain from $B$ to $C$ using lying over, flat going
down and integral going up.  Write $P_i$ for the primes over $p_i$ and
$E_i\subseteq P_i,P_{i+1}$ for the lifted links.  Every $P_i$ is the
kernel of a map $x'_i:C\to A_i$ extending $f_i$.  For clarity, these
maps are local: a nonzero image of an element of the source maximal
ideal in the finite target residue field would, by a monic polynomial
of its inverse, make a unit map to a nonunit.  They are therefore
continuous.  The image of $f_i$ contains $k_i,t_i$ and is compact,
so it is all $A_i$.  The root polynomials split with distinct roots in
$k_i((t_i))$; their roots are in $A_i$.  Normality of $A_i$ then gives
the asserted lift with its exact kernel.  The resulting point of $R/I$
is the original \texttt{OoPt}, with $s\mapsto t_i$, $Y\mapsto0$ and
all roots still distinct.

\emph{Verify the actual patching hypotheses.}
Give $M_N$ its $R$-action through $\rho_N$ and set
\[
 X=C\otimes_BM_0,\qquad
 f=\operatorname{fEl}(1,U)=
                 \prod_j\prod_{i\ne a}(U_{ji}-U_{ja})^2.    \tag{MP8}
\]
If $C$ has rank $r$, then $X$ is the direct sum of $r$ copies of $M_0$
over $\Lambda$, and $YX=0$.  The original equations
$(1+Y_{ji})^{\ell^{a_j(N)}}-1$ lie in
$(Y)^k+\mathfrak m_\Lambda^b$ when $N\ge k+b$.  Thus the original
finite-level freeness gives the exact S5 freeness at every fixed
$k,b$, with the actual possibly varying ranks.

For a fixed coefficient order $b$, finiteness of
$M_0/\mathfrak m_\Lambda^bM_0$ and Nakayama give some $c$ with
$\mathfrak m_B^cM_0\subseteq\mathfrak m_\Lambda^bM_0$.
The limiting determinant and character congruences identify the
original splitting presentations on that quotient.  The actual
Taylor--Wiles maps therefore give the required $R$-linear comparison
maps between $X/\mathfrak m_\Lambda^bX$ and
$M_N/((Y)+\mathfrak m_\Lambda^b)M_N$.
Their composites are multiplication by $f$, since
$\operatorname{fEl}(Q_N,U)\equiv f\pmod{\ell^N}$.

For each relation $a\in\texttt{Rel}$ and each $c$, (MP2) and the
original local-map convergence give, eventually, $a_N\equiv a$
modulo $\mathfrak m_R^c$ whose level-$N$ value is zero or satisfies
the original auxiliary annihilation identity.  The kernel, ordinary
and inertia cases use level-$N$ lifts of the finitely many prescribed
elements; polarization also retains the converging cyclotomic factor.
The root and partial-root cases are exactly TWAut.
If $e_T$ is its single exponent, the difference-of-powers identity gives
\[
 (fa)^{e_T+1}\equiv b_N\pmod{\mathfrak m_R^c},
                   \qquad b_N\in\operatorname{Ann}_RM_N.     \tag{MP9}
\]
For the constructed modules $e_T=n!$.  The source exponent $e_T+1$
is retained, including the zero-exponent edge case.  The quantifiers
are exactly $\forall a\ \forall c\ \text{eventually }N\ \exists b_N$.

At $q_i=\pi^{-1}(P_i)$, $f$ is nonzero because the roots are distinct.
The weight map is onto $A_i$ and sends $s$ to $t_i$.  The completion of
$\Lambda$ at its kernel is regular of dimension $D_{\rm wt}+1$:
$\dim\Lambda=D_{\rm wt}+2$ and its quotient has dimension one.
S5 now supplies its actual support ideals $S_i$ in
\[
 B_i=\widehat{(R/I)_{q_i/I}},\qquad
 \dim B_i/Z\ge 1+D_{\rm wt}+nq=1+D_0                         \tag{MP10}
\]
for every minimal support prime $Z$, together with its exact $Y=0$
fibre and prime-moving property.

\emph{Retain nonsemisimplicity and move nonzero support.}
The limiting group map and $\bar\pi$ give a surjection of the finite
limiting equation algebra onto the original equation algebra after
base change to $C$.  Ordinary, tame and starred equations map to the
originals; the root equation maps to Cayley--Hamilton and the
partial-root equation to zero.  The target is finite and complete,
so it also kills the already closed ideal (MP7).  No saturation or
reduction is used.  Nonsemisimplicity passes to every lifted link:
a nonzero nilpotent ideal survives further field extension, and a
semisimple algebra cannot have a nonsemisimple quotient.  For unrelated
field extensions use a common field over both.  This retains the
source's geometric NonSS predicate even with inseparable residues.

The established B1/B2 comparison gives K1,K2,K3 for each actual OoPt.
GalOut therefore forces $Z\subseteq Q$ whenever $Z$ is stable and
satisfies (MP10), and $Q$ is stable or a lifted nonsemisimple prime.
Moreover finite-module support and $\pi$ surjective give
\[
 \operatorname{Ann}_RX\subseteq\pi^{-1}(P)
                   \quad\Longleftrightarrow\quad J\subseteq P\cap B.
                                                               \tag{MP11}
\]
Lift the initial $q_0$ to $B_0$.  The S5 fibre identity and (MP11)
give a nonzero support prime.  Choose a minimal support prime below
it; (MP1) preserves stability under generalization and (MP10) gives
its dimension.  GalOut puts it below a prime over the first link.
S5 moves it to the next completion with the same contraction.
Repeat through the whole finite chain.  At $q_*$ apply GalOut once
more, then the fibre identity and (MP11), to obtain $J\subseteq q_*$.
This proves marked Core, without assuming that support is nonempty
before the given initial prime is used.

\emph{Remove the marking and coefficient restrictions.}
Choose the unramified finite extension $O'/O$ with residue field of
the common cardinality $|k_i|$, and put
\[
                         B'=(O'\otimes_OB)[[s]].             \tag{MP12}
\]
The original IsCO condition makes $O/\mathfrak m_O\to B/\mathfrak m_B$
an isomorphism.  Thus $(O'\otimes_OB)/\mathfrak m_B(O'\otimes_OB)$
is the field $O'/\mathfrak m_{O'}$, proving that this finite complete
algebra is local.  Base change every original module and map, and
extend in this one formal variable.  Freeness, diamond quotients, compositions and the
single exponent are unchanged.  First lift the entire prime chain
through the finite flat coefficient extension.  The chosen monic
generator polynomial of $O'$ splits in each $k_i$; each lifted prime
thus gives a map into the same $A_i$, with $O'\twoheadrightarrow k_i$.
Send $s$ to $t_i$ and extend the linking, initial and terminal primes
as power-series ideals.  They remain prime and lie below the required
curve kernels.  The finite base-changed equation algebra kills the
new closed relations continuously, so the same quotient argument
retains NonSS.  Apply Core and contract to $B$.  This is CoreBC.

Finally quotient each open residual kernel $H$ by the kernel of its
maximal pro-$\ell$ quotient.  The finite $H^1(H,\mathbb F_\ell)$
bound gives finite residual-adjoint $H^1$ by restriction and finitely
many coset values, so the continuous Cayley--Hamilton algebra is
finite even before this quotient.  Since all residual characters
are trivial on $H$, its image lies in a pro-$\ell$ group: modulo
the coefficient maximal ideal it is in the nilpotent radical, and
higher congruence kernels are pro-$\ell$.  Thus every determinant,
point and auxiliary action factors through the quotient.  A finite
Selmer cocycle together with its coefficient action has values in
a finite semidirect $\ell$-group on $H$ and factors as well.  This
identifies the original cocycles, local witnesses and boundaries,
so B1/B2 and the odd families descend without loss.  The finite
equation algebras, their closed ideals and (MP1) also descend exactly.
CoreBC applied to this tower gives Raw for the original data.
\end{proof}

This closes the actual marked-model and chain-propagation supplier.
The initial stable prime below a prescribed coefficient curve remains
an explicit input.  The subsequent global component argument and
its uniform integral Selmer budget are not consequences of a nonzero
characteristic-zero automorphic point alone.
